\chapter{A Concrete Naming Certificate for $\EffectiveCauchyModulusUp$}
\label{ch:concrete-instances-effective-cauchy-modulus-namecert}
\origin{human}

Following Bishop and Bridges, \emph{Constructive Analysis}, the $\EffectiveCauchyModulusUp$ packet records the finite precision-to-tail schedule that lets a Cauchy name serve L10 Real and Completion consumers without importing ambient completeness. It sits over the generic modulus schedule of \autoref{ch:concrete-instances-cauchymodulus-namecert}, the fast threshold discipline of \autoref{ch:concrete-instances-fast-cauchy-modulus-namecert}, the request-to-tail realizer of \autoref{ch:concrete-instances-regular-cauchy-modulus-realizer-namecert}, the finite-window interface of \autoref{ch:concrete-instances-streamname-namecert}, the regular rational route of \autoref{ch:concrete-instances-regseqrat-namecert}, and the terminal $\RealUp$ seal. The packet supplies one displayed schedule row before any $\RegSeqRatUp$ or $\RealUp$ consumer reads a seal.

\begin{definition}[Effective Cauchy modulus carrier]
\label{def:effective-cauchy-modulus-carrier}
An $\EffectiveCauchyModulusUp$ carrier is a finite $\BHist$ packet
$$
X=(S,D,M,T,W,Q,E,H,C,P,N)
$$
with the following displayed rows:
\begin{enumerate}
\item $S$ is the $\StreamNameUp$ source row whose observations are opened only through finite windows.
\item $D$ is the dyadic precision ledger consumed by the tail estimates and regular rational readbacks.
\item $M$ is the effective precision-to-tail modulus row, assigning each displayed precision request a finite tail index.
\item $T$ is the tail-index comparison row showing that the selected index is stable under stronger displayed precision requests.
\item $W$ is the finite tail-window row opened from $S$ at the index selected by $M$.
\item $Q$ is the $\RegSeqRatUp$ readback row obtained from $W$ under the dyadic ledger $D$.
\item $E$ is the $\RealUp$ handoff row, exposed only after $S,D,M,T,W,Q$ have been displayed.
\item $H$ records componentwise $\hsame$ transport for the stream, dyadic ledger, effective modulus, tail comparison, window, readback, handoff, replay, provenance, and local naming rows.
\item $C$ records displayed $\Cont$ replay from a precision request through $S,D,M,T,W,Q,E_H,H,P,N$.
\item $P$ is the $\Pkg$ provenance over $\EffectiveCauchyModulusUp$, $\CauchyModulusUp$, $\FastCauchyModulusUp$, $\RegularCauchyModulusRealizerUp$, $\StreamNameUp$, $\RegSeqRatUp$, $\RealUp$, $\BHist$, $\hsame$, $\Cont$, and $\NameCert$.
\item $N$ is the local $\NameCert_{\EffectiveCauchyModulusUp}$ row.
\end{enumerate}
The classifier compares accepted packets only through $S,D,M,T,W,Q,E,H,C,P,N$. It has no coordinate for quotient limits, countable choice, host convergence, trichotomy, selected representatives, or ambient $\RealUp$ completeness.
\end{definition}

\begin{theorem}[Effective Cauchy modulus NameCert obligations]
\label{thm:effective-cauchy-modulus-namecert-obligations}
For every accepted $\EffectiveCauchyModulusUp$ carrier
$$
X=(S,D,M,T,W,Q,E,H,C,P,N),
$$
the local $\NameCert_{\EffectiveCauchyModulusUp}$ surface consists exactly of the following finite obligations:
\begin{enumerate}
\item carrier exposure for the stream source, dyadic precision ledger, effective modulus row, tail comparison row, finite window, regular readback, Real handoff, transport, replay, provenance, and local naming rows;
\item schedule classification, requiring $M$ to be the only precision-to-tail row read by a downstream $\RegSeqRatUp$ or $\RealUp$ consumer;
\item monotone tail-index stability, requiring $T$ to certify that stronger displayed precision requests never require a weaker tail boundary than the row already exposed by $M$;
\item $\RegSeqRatUp$ and $\RealUp$ seal compatibility, requiring $W$ and $Q$ to be read before the handoff row $E$ exposes a terminal Real seal;
\item ledger refusal for quotient limits, countable choice, host convergence, trichotomy, selected representatives, and ambient completeness as local certificate data.
\end{enumerate}
No effective Cauchy modulus obligation is stored outside these rows.
\end{theorem}
\begin{proof}
Project the accepted carrier. The analytic coordinates are exactly the stream source, dyadic ledger, effective modulus, tail comparison, finite window, regular readback, and Real handoff rows.

The schedule classifier is the row $M$. Since every consumer reaches a tail window only after reading $M$, no hidden convergence selector or external completeness principle can choose the boundary.

The monotonicity information is stored in $T$. It relates stronger displayed precision requests to tail indices already chosen by $M$, so transport through $H$ and replay through $C$ preserve the same schedule.

Finally, the readback and seal are downstream of the finite window. The rows $Q$ and $E$ can expose regular rational and Real-facing data only after $S,D,M,T,W$ have been displayed, and the refused objects have no carrier coordinate.
\end{proof}

\begin{theorem}[Effective Cauchy modulus seal route]
\label{thm:effective-cauchy-modulus-seal-route}
Every $\RegSeqRatUp$ or $\RealUp$ consumer of an accepted $\EffectiveCauchyModulusUp$ packet factors through
$$
S,D,M,T \longmapsto W \longmapsto Q \longmapsto E \longmapsto H,C,P,N .
$$
The route exports a displayed finite tail schedule and does not export quotient limits, countable choice, host convergence, trichotomy, selected representatives, or ambient $\RealUp$ completeness.
\end{theorem}
\begin{proof}
Read the carrier coordinates in order. A precision request first sees the stream and dyadic rows $S,D$, then the effective modulus and tail-stability rows $M,T$.

Only after those rows are displayed can the finite window $W$ be opened. The regular rational readback $Q$ and terminal handoff $E_H$ are downstream of that finite window.

The structural rows $H,C,P,N$ transport, replay, source, and name the same route. Since the carrier contains no coordinate for the refused principles, the exported object remains the finite schedule-to-window packet.
\end{proof}

\closureat{EffectiveCauchyModulusUp}{seedStr}

\begin{closurestatus}{\EffectiveCauchyModulusUp}
  \constructivestory{An $\EffectiveCauchyModulusUp$ packet is generated from a StreamName source row, a dyadic precision ledger, an effective precision-to-tail modulus row, a monotone tail-index comparison row, a finite tail-window row, a RegSeqRat readback, a RealUp handoff, componentwise hsame transport, displayed Cont replay, Pkg provenance, and a local NameCert row.}
  \theoryclosure{\seedClosure}
  \scopeclosed{\autoref{def:effective-cauchy-modulus-carrier}, \autoref{thm:effective-cauchy-modulus-namecert-obligations}, and \autoref{thm:effective-cauchy-modulus-seal-route} seed the finite schedule object consumed by L10 StreamName, RegSeqRat, and Real routes.}
  \formalstatus{\unformalizedV}
  \bridgestatus{none}
  \notclaimed{This chapter does not close RealUp completeness, quotient limits, countable choice, host convergence, trichotomy, selected representatives, Lean verification, or bridge checking.}
  \upgradepath{Obligation closure requires checked carrier exposure, schedule classification, monotone tail-index stability, RegSeqRat and RealUp seal compatibility, non-escape, transport, replay, provenance, and local naming for BEDC.Derived.EffectiveCauchyModulusUp.}
\end{closurestatus}
